\documentclass[11pt,a4paper]{article}

% ============================================================
%  QIMMA -- COURS : Les suites numériques (2BAC)
%  Basé sur le template cours.tex (charte Qimma)
% ============================================================

\usepackage[french]{babel}
\usepackage[utf8]{inputenc}
\usepackage[T1]{fontenc}
\usepackage{lmodern}
\usepackage[a4paper,margin=2cm,top=2.3cm,bottom=2.6cm]{geometry}
\usepackage{amsmath,amssymb}
\usepackage{xcolor}
\usepackage{graphicx}
\usepackage{tikz}
\usetikzlibrary{shadows,calc}
\usepackage[most]{tcolorbox}
\usepackage{titlesec}
\usepackage{fancyhdr}
\usepackage{enumitem}
\usepackage{pifont}
\usepackage{array}
\usepackage{colortbl}
\usepackage{hyperref}

% ------------------- CHARTE QIMMA -------------------
\definecolor{qteal}{HTML}{0d9488}
\definecolor{qtealdark}{HTML}{134e4a}
\definecolor{qamber}{HTML}{fbbf24}
\definecolor{qambertext}{HTML}{92400e}
\definecolor{ink}{HTML}{1f2937}
\definecolor{inkgray}{HTML}{6b7280}
\definecolor{tealpale}{HTML}{e6f5f3}
\colorlet{colDef}{qteal}
\definecolor{colProp}{HTML}{0f766e}
\definecolor{colEx}{HTML}{b45309}
\definecolor{colRem}{HTML}{9d174d}
\color{ink}

\graphicspath{{assets/}}
\newcommand{\logofile}{assets/qimma-logo.pdf}

% ------------------- METADONNEES -------------------
\newcommand{\matiere}{Mathématiques}
\newcommand{\niveau}{2\up{ème} Bac -- PC, SVT, SM}
\newcommand{\chapitretitre}{Les suites numériques}

% ------------------- EN-TETE / PIED DE PAGE -------------------
\pagestyle{fancy}
\fancyhf{}
\renewcommand{\headrulewidth}{0pt}
\renewcommand{\footrulewidth}{0.5pt}
\fancyfoot[L]{\raisebox{-0.15cm}{\includegraphics[height=0.35cm]{\logofile}}~\small\sffamily\color{inkgray}Qimma}
\fancyfoot[C]{\small\sffamily\color{inkgray}\matiere\ -- \niveau}
\fancyfoot[R]{\small\sffamily\color{qtealdark}\thepage}

% ------------------- TITRES DE SECTION -------------------
\titleformat{\section}
  {\normalfont\sffamily\Large\bfseries\color{qtealdark}}
  {\tikz[baseline=-3.2pt]{\node[circle,fill=qteal,text=white,inner sep=0pt,minimum size=8mm]{\sffamily\bfseries\thesection};}}
  {10pt}{}
  [\vspace{-2mm}{\color{qamber}\titlerule[1.6pt]}\vspace{2mm}]
\titlespacing*{\section}{0pt}{16pt}{10pt}

\titleformat{\subsection}
  {\normalfont\sffamily\large\bfseries\color{qtealdark}}
  {\color{qteal}\ding{212}}{6pt}{}
\titlespacing*{\subsection}{0pt}{12pt}{6pt}

% ------------------- BOITES PEDAGOGIQUES -------------------
\newtcolorbox{fiche}[3][]{
  enhanced, breakable,
  colback=white, colframe=#2,
  boxrule=1pt, arc=2mm,
  top=7mm, left=4mm, right=4mm, bottom=3mm,
  drop shadow={black!25, shadow xshift=1mm, shadow yshift=-1mm},
  attach boxed title to top left={xshift=4mm, yshift=-3mm, yshifttext=-1mm},
  boxed title style={colback=#2, arc=1.5mm, boxrule=0pt, left=2mm, right=2mm},
  coltitle=white, fonttitle=\sffamily\bfseries,
  title=#3, #1
}
\newenvironment{definition}[1][Définition]
  {\begin{fiche}{colDef}{\ding{110}~~#1}}{\end{fiche}}
\newenvironment{propriete}[1][Propriété]
  {\begin{fiche}{colProp}{\ding{72}~~#1}}{\end{fiche}}
\newenvironment{exemple}[1][Exemple]
  {\begin{fiche}{colEx}{\ding{217}~~#1}}{\end{fiche}}
\newenvironment{remarque}[1][Remarque]
  {\begin{fiche}{colRem}{\ding{43}~~#1}}{\end{fiche}}

\hypersetup{colorlinks=true, linkcolor=qtealdark, urlcolor=qtealdark}

% ============================================================
\begin{document}

% ------------------- BANNIERE DE TITRE -------------------
\begin{tcolorbox}[
  enhanced, boxrule=0pt, arc=4mm,
  interior style={left color=qtealdark, right color=qteal},
  top=8mm, bottom=8mm, left=10mm, right=32mm,
  drop shadow={black!35, shadow xshift=1.5mm, shadow yshift=-1.5mm},
  before skip=0pt, after skip=9mm,
  overlay={
    \node[anchor=north west] at ([xshift=6mm,yshift=-6mm]frame.north west)
      {\includegraphics[height=1.25cm]{\logofile}};
    \node[anchor=north east, fill=qamber, text=qambertext, rounded corners=2pt,
          inner sep=4pt, font=\sffamily\bfseries\small] at ([xshift=-6mm,yshift=-6mm]frame.north east) {COURS};
  }
]
\hspace{1.6cm}\centering
{\sffamily\bfseries\color{white}\fontsize{23}{27}\selectfont \chapitretitre}\\[3mm]
{\sffamily\color{white!90}\large \matiere\ \textbullet\ \niveau}
\end{tcolorbox}

% ------------------- OBJECTIFS -------------------
\begin{tcolorbox}[
  enhanced, colback=tealpale, colframe=qteal, boxrule=0.8pt, arc=2mm,
  left=5mm, right=5mm, top=3mm, bottom=3mm,
  title={\sffamily\bfseries\color{qtealdark}\ding{51}~~Objectifs du chapitre},
  coltitle=qtealdark, colbacktitle=tealpale, boxrule=0.8pt,
  attach title to upper={\par\vspace{1mm}}
]
\begin{itemize}[leftmargin=6mm, label=\ding{212}, itemsep=1mm]
  \item Reconnaître une suite et la définir de façon explicite ou par récurrence.
  \item Identifier et exploiter une suite arithmétique ou géométrique (terme général, somme).
  \item Étudier le sens de variation et le caractère borné d'une suite.
  \item Calculer une limite à l'aide des limites de référence et des opérations sur les limites.
  \item Utiliser les théorèmes de convergence (monotonie, comparaison, encadrement).
  \item Traiter une suite arithmético-géométrique du type $u_{n+1} = a\,u_n + b$.
\end{itemize}
\end{tcolorbox}

\vspace{4mm}

% ------------------- SECTION 1 : GENERALITES -------------------
\section{Généralités sur les suites}

\begin{definition}
Une \textbf{suite numérique} est une fonction qui associe, à chaque entier naturel $n$, un nombre réel noté $u_n$. On note la suite $(u_n)$ pour la distinguer de son \textbf{terme général} $u_n$, qui est un nombre. Le premier terme (souvent $u_0$, parfois $u_1$) est appelé \textbf{terme initial}.
\end{definition}

Il existe deux grandes façons de définir une suite, et savoir les distinguer est essentiel car les méthodes d'étude ne sont pas les mêmes :

\begin{enumerate}[leftmargin=8mm, itemsep=1mm]
  \item \textbf{Définition explicite} : une formule directe $u_n = f(n)$ calcule n'importe quel terme sans connaître les précédents. Exemple : $u_n = 2n + 3$.
  \item \textbf{Définition par récurrence} : on donne le premier terme et une relation qui calcule chaque terme à partir du précédent, $u_{n+1} = f(u_n)$. Exemple : $u_0 = 1$ et $u_{n+1} = 2u_n + 1$.
\end{enumerate}

\begin{remarque}[Astuce]
Si tu peux calculer $u_{100}$ directement en remplaçant $n$ par $100$, la définition est \emph{explicite}. Si tu as besoin de $u_{99}$ pour obtenir $u_{100}$, c'est une \emph{récurrence}.
\end{remarque}

% ------------------- SECTION 2 : ARITHMETIQUES -------------------
\section{Suites arithmétiques}

\begin{definition}
Une suite $(u_n)$ est \textbf{arithmétique} lorsqu'on passe d'un terme au suivant en ajoutant toujours le même réel $r$, appelé la \textbf{raison} :
\[
  u_{n+1} = u_n + r \qquad \text{pour tout } n \in \mathbb{N}.
\]
\end{definition}

\begin{propriete}[Terme général et somme]
Pour une suite arithmétique de raison $r$ :
\[
  u_n = u_0 + n\,r
  \qquad\text{et plus généralement}\qquad
  u_n = u_p + (n - p)\,r.
\]
La somme de termes consécutifs vaut :
\[
  S = (\text{nombre de termes}) \times \frac{\text{premier terme} + \text{dernier terme}}{2}.
\]
En particulier : $1 + 2 + \cdots + n = \dfrac{n(n+1)}{2}$.
\end{propriete}

\begin{exemple}[Méthode attendue au bac]
Soit $u_n = 5n - 2$. Montrons que $(u_n)$ est arithmétique :
\[
  u_{n+1} - u_n = \bigl[5(n+1) - 2\bigr] - \bigl[5n - 2\bigr] = 5n + 5 - 2 - 5n + 2 = 5.
\]
La différence est constante : $(u_n)$ est arithmétique de raison $r = 5$ et de premier terme $u_0 = -2$.
\end{exemple}

% ------------------- SECTION 3 : GEOMETRIQUES -------------------
\section{Suites géométriques}

\begin{definition}
Une suite $(u_n)$ est \textbf{géométrique} lorsqu'on passe d'un terme au suivant en multipliant toujours par le même réel $q$, appelé la \textbf{raison} :
\[
  u_{n+1} = q \cdot u_n \qquad \text{pour tout } n \in \mathbb{N}.
\]
\end{definition}

\begin{propriete}[Terme général et somme]
Pour une suite géométrique de raison $q$ :
\[
  u_n = u_0 \cdot q^n
  \qquad\text{et plus généralement}\qquad
  u_n = u_p \cdot q^{\,n-p}.
\]
Somme de termes consécutifs (pour $q \neq 1$) :
\[
  S = \text{premier terme} \times \frac{1 - q^{\text{nombre de termes}}}{1 - q},
  \qquad\text{en particulier}\quad
  1 + q + q^2 + \cdots + q^n = \frac{1 - q^{n+1}}{1 - q}.
\]
\end{propriete}

\begin{remarque}[Ne confonds pas les deux modèles]
Arithmétique = on \textbf{ajoute} $r$ (croissance « en ligne droite »). Géométrique = on \textbf{multiplie} par $q$ (croissance ou décroissance « accélérée »). Pour prouver la nature d'une suite : calcule $u_{n+1} - u_n$ (arithmétique) ou $\dfrac{u_{n+1}}{u_n}$ (géométrique) et montre que le résultat ne dépend pas de $n$.
\end{remarque}

% ------------------- SECTION 4 : VARIATION ET BORNES -------------------
\section{Sens de variation et suites bornées}

\subsection{Sens de variation}

Trois méthodes permettent d'étudier la monotonie d'une suite ; on choisit selon sa forme :

\begin{enumerate}[leftmargin=8mm, itemsep=1mm]
  \item Étudier le \textbf{signe de $u_{n+1} - u_n$} : positif pour tout $n$ $\Rightarrow$ croissante ; négatif $\Rightarrow$ décroissante. \emph{Méthode universelle, à privilégier.}
  \item Comparer le \textbf{rapport $\dfrac{u_{n+1}}{u_n}$ à $1$} : utile \emph{uniquement} si tous les termes sont strictement positifs.
  \item Utiliser la \textbf{fonction associée} : si $u_n = f(n)$ avec $f$ monotone sur $[0\,;+\infty[$, la suite a le même sens de variation que $f$.
\end{enumerate}

\subsection{Suites majorées, minorées, bornées}

\begin{definition}
La suite $(u_n)$ est \textbf{majorée} s'il existe un réel $M$ tel que $u_n \leq M$ pour tout $n$ ; \textbf{minorée} s'il existe un réel $m$ tel que $u_n \geq m$ pour tout $n$ ; \textbf{bornée} si elle est à la fois majorée et minorée : $m \leq u_n \leq M$.
\end{definition}

\begin{remarque}[Erreur classique]
Un majorant n'est pas forcément atteint : dire « $u_n \leq 3$ » ne signifie pas qu'un terme vaut $3$. Et il ne suffit pas de vérifier sur les premiers termes : il faut démontrer l'inégalité pour \emph{tout} $n$.
\end{remarque}

% ------------------- SECTION 5 : LIMITES -------------------
\section{Limite d'une suite et convergence}

\begin{definition}
Si les termes $u_n$ se stabilisent autour d'un réel $\ell$ quand $n \to +\infty$, on dit que la suite \textbf{converge} vers $\ell$ et on écrit $\lim\limits_{n \to +\infty} u_n = \ell$. Sinon, la suite \textbf{diverge} (elle tend vers $+\infty$, vers $-\infty$, ou n'a pas de limite).
\end{definition}

\begin{propriete}[Limites de référence]
\[
  \frac{1}{n} \to 0, \qquad
  \frac{1}{n^2} \to 0, \qquad
  \frac{1}{\sqrt{n}} \to 0, \qquad
  \sqrt{n} \to +\infty, \qquad
  n^k \to +\infty \ (k > 0).
\]
Pour la suite géométrique $(q^n)$ :
\begin{itemize}[leftmargin=8mm, itemsep=0.5mm]
  \item si $-1 < q < 1$, alors $q^n \to 0$ ;
  \item si $q > 1$, alors $q^n \to +\infty$ ;
  \item si $q = 1$, alors $q^n = 1$ pour tout $n$ ;
  \item si $q \leq -1$, la suite $(q^n)$ n'a pas de limite.
\end{itemize}
\end{propriete}

\begin{remarque}[Les quatre formes indéterminées]
« $\infty - \infty$ », « $0 \times \infty$ », « $\dfrac{\infty}{\infty}$ » et « $\dfrac{0}{0}$ ». Face à une forme indéterminée, on ne conclut jamais directement : il faut transformer l'expression (factoriser par le terme dominant, simplifier, utiliser l'expression conjuguée\ldots).
\end{remarque}

\begin{exemple}[Lever une indétermination $\infty/\infty$]
Cherchons la limite de $u_n = \dfrac{3n^2 + 2n}{n^2 + 5}$. Numérateur et dénominateur tendent vers $+\infty$ : forme indéterminée. On factorise par le terme dominant $n^2$ :
\[
  u_n = \frac{n^2\left(3 + \frac{2}{n}\right)}{n^2\left(1 + \frac{5}{n^2}\right)}
      = \frac{3 + \frac{2}{n}}{1 + \frac{5}{n^2}}
  \xrightarrow[n \to +\infty]{} \frac{3 + 0}{1 + 0} = 3.
\]
La suite converge vers $3$.
\end{exemple}

% ------------------- SECTION 6 : THEOREMES -------------------
\section{Théorèmes de convergence}

\begin{propriete}[Convergence monotone]
Toute suite \textbf{croissante et majorée} converge. Toute suite \textbf{décroissante et minorée} converge.
\end{propriete}

\begin{remarque}
Ce théorème prouve que la limite \emph{existe}, mais ne donne pas sa valeur.
\end{remarque}

\begin{propriete}[Théorème des gendarmes (encadrement)]
Si $v_n \leq u_n \leq w_n$ à partir d'un certain rang, et si $(v_n)$ et $(w_n)$ tendent vers la même limite $\ell$, alors $(u_n)$ tend aussi vers $\ell$. \emph{Idéal pour les suites contenant $\cos(n)$ ou $\sin(n)$, que l'on encadre entre $-1$ et $1$.}
\end{propriete}

\begin{propriete}[Théorème de comparaison]
Si $u_n \geq v_n$ à partir d'un certain rang et si $v_n \to +\infty$, alors $u_n \to +\infty$. De même, si $u_n \leq v_n$ et $v_n \to -\infty$, alors $u_n \to -\infty$.
\end{propriete}

% ------------------- SECTION 7 : ARITHMETICO-GEOMETRIQUES -------------------
\section{Suites arithmético-géométriques}

C'est le type de suite le plus fréquent dans les problèmes de bac : $u_{n+1} = a\,u_n + b$ (avec $a \neq 1$ et $b \neq 0$). Elle n'est ni arithmétique ni géométrique, mais une méthode standard la ramène à une suite géométrique.

\begin{propriete}[Méthode en 4 étapes]
\begin{enumerate}[leftmargin=8mm, itemsep=1mm]
  \item \textbf{Point fixe} : résoudre $\ell = a\,\ell + b$. C'est la limite éventuelle de la suite.
  \item \textbf{Suite auxiliaire} : poser $v_n = u_n - \ell$ et montrer qu'elle est géométrique ; on trouve $v_{n+1} = a\,v_n$, de raison $a$.
  \item \textbf{Terme général} : $v_n = v_0 \cdot a^n$, puis revenir à $u_n = v_n + \ell$.
  \item \textbf{Limite} : si $-1 < a < 1$, alors $a^n \to 0$ donc $u_n \to \ell$.
\end{enumerate}
\end{propriete}

\begin{exemple}[Étude complète]
Soit $u_0 = 5$ et $u_{n+1} = \frac{1}{2}u_n + 1$.
\begin{itemize}[leftmargin=6mm, itemsep=0.5mm]
  \item Point fixe : $\ell = \frac{1}{2}\ell + 1$ donne $\frac{1}{2}\ell = 1$, soit $\ell = 2$.
  \item On pose $v_n = u_n - 2$. Alors $v_{n+1} = u_{n+1} - 2 = \frac{1}{2}u_n + 1 - 2 = \frac{1}{2}(u_n - 2) = \frac{1}{2}v_n$.
  \item $(v_n)$ est géométrique de raison $\frac{1}{2}$, avec $v_0 = u_0 - 2 = 3$, d'où $v_n = 3 \cdot \left(\frac{1}{2}\right)^n$.
  \item Donc $u_n = 2 + 3 \cdot \left(\frac{1}{2}\right)^n$. Comme $-1 < \frac{1}{2} < 1$, on a $\left(\frac{1}{2}\right)^n \to 0$.
\end{itemize}
Conclusion : $\lim u_n = 2$, la suite converge vers $2$.
\end{exemple}

\begin{remarque}[Astuce mémo]
La suite auxiliaire a \emph{toujours} pour raison le coefficient $a$ de la relation $u_{n+1} = a\,u_n + b$, et le nombre à retrancher est \emph{toujours} le point fixe $\ell$.
\end{remarque}

% ------------------- SECTION 8 : METHODES -------------------
\section{Méthodes à connaître}

\begin{center}
\renewcommand{\arraystretch}{1.45}
\sffamily\small
\begin{tabular}{>{\raggedright\arraybackslash}p{6.6cm} >{\raggedright\arraybackslash}p{9.4cm}}
\rowcolor{qtealdark} \color{white}\bfseries Ce qu'on te demande & \color{white}\bfseries Ce que tu fais \\
\rowcolor{tealpale} Montrer qu'une suite est arithmétique & Calculer $u_{n+1} - u_n$ et montrer que c'est une constante $r$. \\
Montrer qu'une suite est géométrique & Calculer $u_{n+1}/u_n$ et montrer que c'est une constante $q$. \\
\rowcolor{tealpale} Étudier le sens de variation & Étudier le signe de $u_{n+1} - u_n$ (méthode par défaut). \\
Calculer une limite & Limites de référence + opérations ; lever toute indétermination. \\
\rowcolor{tealpale} Prouver la convergence sans calculer & Montrer que la suite est monotone et bornée. \\
Limite avec $\cos(n)$ ou $\sin(n)$ & Encadrer, puis appliquer le théorème des gendarmes. \\
\rowcolor{tealpale} Suite $u_{n+1} = a\,u_n + b$ & Point fixe $\ell$, suite auxiliaire $v_n = u_n - \ell$ (géométrique de raison $a$). \\
\end{tabular}
\end{center}

\begin{remarque}[Les erreurs qui coûtent des points au bac]
\begin{itemize}[leftmargin=6mm, itemsep=0.5mm]
  \item Conclure « la limite est $\ell$ » à partir de $\ell = a\ell + b$ \emph{sans} avoir prouvé que la suite converge.
  \item Utiliser le quotient $u_{n+1}/u_n$ pour la variation alors que les termes ne sont pas tous positifs.
  \item Appliquer la formule de somme géométrique quand $q = 1$ (division par zéro !) : si $q = 1$, la somme vaut $n \times$ premier terme.
  \item Écrire « $\infty - \infty = 0$ » ou « $\infty/\infty = 1$ » : ce sont des formes indéterminées, jamais un résultat.
  \item Confondre $u_n$ (un nombre) et $(u_n)$ (la suite entière) dans la rédaction.
\end{itemize}
\end{remarque}

% ------------------- RESUME -------------------
\vspace{4mm}
\begin{tcolorbox}[
  enhanced, boxrule=0pt, arc=2mm,
  interior style={left color=qtealdark, right color=qteal},
  left=6mm, right=6mm, top=4mm, bottom=4mm,
  fontupper=\color{white}\sffamily,
  title={\sffamily\bfseries\color{white}\ding{72}~~À retenir},
  colbacktitle=qtealdark, coltitle=white, boxrule=0pt
]
\begin{itemize}[leftmargin=6mm, label={\color{qamber}\ding{212}}, itemsep=1mm]
  \item Arithmétique : $u_{n+1} = u_n + r$, terme général $u_n = u_0 + n\,r$.
  \item Géométrique : $u_{n+1} = q\,u_n$, terme général $u_n = u_0\,q^n$.
  \item Sens de variation : signe de $u_{n+1} - u_n$ (méthode par défaut).
  \item Convergence garantie : suite croissante majorée (ou décroissante minorée).
  \item Limite de $q^n$ : $0$ si $|q| < 1$ ; $+\infty$ si $q > 1$.
  \item Formes indéterminées : $\infty - \infty$, $0 \times \infty$, $\infty/\infty$, $0/0$.
  \item Arithmético-géométrique $u_{n+1} = a\,u_n + b$ : point fixe $\ell$, puis $v_n = u_n - \ell$ géométrique de raison $a$.
\end{itemize}
\end{tcolorbox}

\end{document}
